\section{Main results}
\label{sect:main}

We start with a~new notion of ``mystical equivalence'' of group actions, which is crucial for what follows.

\begin{Definition}
Let $ \mathop{\triangleright}\nolimits\colon G \times \mathcal V\to \mathcal V$,
$\mathop{\triangleright}\nolimits\nolimits'\colon G' \times \mathcal V\to \mathcal V $ be faithful actions of f\/inite
groups~$G$, respectively $G'$, on a~complex vector space~$\mathcal V$.
We say that the actions $\mathop{\triangleright}\nolimits$ and $\mathop{\triangleright}\nolimits'$ are
\textit{mystically equivalent}, if
\begin{gather*}
\rho(e_G)=\rho'(e_{G'}),
\end{gather*}
where $\rho\colon \mathbb C G \to \End_\mathbb C \mathcal V$ and $\rho'\colon \mathbb C G'\to \End_\mathbb C\mathcal V$
are the algebra homomorphisms def\/ined by the actions, and $e_G$ denotes the element $\sum\limits_{g\in G}g$ of $\mathbb C G$.
\end{Definition}

Mystical equivalence of the actions of~$G$ and $G'$ is a~strengthening of the condition that the respective spaces
$\mathcal V^G$ and $\mathcal V^{G'}$ of invariants are equal, due to the following obvious result.
\begin{Lemma}%\label{lemma:projection}
If a~$G$-action and a~$G'$-action on $\mathcal V$ are mystically equivalent, then $\mathcal V^G=\mathcal V^{G'}$.
\end{Lemma}

We will use the lemma in the situation where $\mathcal V = S(V)$ where $V$ is a~vector space over $\mathbb C$ with
a~chosen basis $\{x_1,\ldots,x_n\}$.
Throughout the paper, $n\ge 2$.
Denote by $\mathbb G_n$ the group of monomial matrices on $V$, that is, matrices in $\mathrm{GL}_n(\mathbb C)$ with
exactly $n$ non-zero entries.
In other words,
\begin{gather*}
\mathbb G_n=(\mathbb C^\times)^n\rtimes \mathbb S_n.
\end{gather*}
Here $(\mathbb C^\times)^n$ is naturally identif\/ied with the group of diagonal matrices in $\mathrm{GL}_n(\mathbb C)$
and acts on~$V$ by scaling the basis $\{x_1,\ldots,x_n\}$.
The symmetric group $\mathbb S_n$ is identif\/ied with the group of permutation matrices and acts on~$V$ by permuting the
same basis.
Note that $\mathbb G_n$ is the normalizer of the torus $(\mathbb C^\times)^n$ in $\mathrm{GL}_n(\mathbb C)$.
In particular, $\mathbb S_n$ acts on $(\mathbb C^\times)^n$ by conjugation.
When we write $tw\in \mathbb G_n$, we will imply that $t\in (\mathbb C^\times)^n$ and $w\in \mathbb S_n$; every element
of $\mathbb G_n$ can be uniquely written in this way.

Clearly, $\mathbb G_n$ is generated by $s_1,\ldots,s_{n-1}$ and $t_j^{(\zeta)}$, $1\le j\le n$, $\zeta\in \mathbb
C^\times$, where
\begin{itemize}\itemsep=0pt
\item
$s_i\in \mathbb S_n$ is the permutation of $\{x_1,\ldots,x_n\}$ which swaps $x_i$ and $x_{i+1}$;
\item
$t_j^{(\zeta)}\in (\mathbb C^\times)^n$ maps $x_k$ to $\zeta^{\delta_{jk}} x_k$.
\end{itemize}
We f\/ind it very convenient to use the linear character
\begin{gather*}
\det \colon \ \mathbb G_n \to \mathbb C^\times
\end{gather*}
of $\mathbb G_n$, which is just the restriction of the determinant character of $\mathrm{GL}_n(\mathbb C)$ to $\mathbb
G_n$.
In particular,
\begin{gather*}
\det s_i = -1,
\qquad
\det w\in\{\pm1\}\text{ for }w\in \mathbb S_n,
\qquad
\det\left(t_1^{(\zeta_1)} t_2^{(\zeta_2)} \cdots t_n^{(\zeta_n)}\right)=\zeta_1\zeta_2\cdots \zeta_n.
\end{gather*}

{\sloppy Next, we introduce \textit{two} dif\/ferent faithful actions of $\mathbb G_n$ on $S(V)$ using the natural basis
$\big\{x_1^{k_1} \cdots x_n^{k_n} :  k_1,\ldots,k_n\in \mathbb Z_{\ge0}\big\}$ of~$S(V)$.
(At the moment, we are not using any multiplication on $S(V)$.)

}

\begin{Proposition}\label{prop:automorphism}\quad
\begin{enumerate}\itemsep=0pt
\item[$(a)$] There exist $($unique$)$ faithful actions $\mathop{\triangleright}\nolimits_+$, $\mathop{\triangleright}\nolimits_-$ of
$\mathbb G_n$ on $S(V)$ such that
\begin{gather*}
s_i \mathop{\triangleright}\nolimits_\pm x_1^{k_1}\cdots x_n^{k_n} = (\pm 1)^{k_ik_{i+1}} x_1^{k_1}\cdots
x_i^{k_{i+1}}x_{i+1}^{k_i}\cdots x_n^{k_n},
\\
t_j^{(\zeta)}\mathop{\triangleright}\nolimits_\pm x_1^{k_1}\cdots x_n^{k_n} = \zeta^{k_j} x_1^{k_1}\cdots
x_n^{k_n}
\end{gather*}
for any $i=1,\ldots,n-1$, $j=1,\ldots,n$, $\zeta\in \mathbb C^\times$.
Both actions extend the defining action of~$\mathbb G_n$ on~$V$.

\item[$(b)$] The action $\mathop{\triangleright}\nolimits_+$ of $\mathbb G_n$ on $S(V)$ is compatible with the natural
commutative multiplication on~$S(V)$, in the sense that~$\mathbb G_n$ acts by automorphisms of the algebra~$S(V)$.

\item[$(c)$] The action $\mathop{\triangleright}\nolimits_-$ of $\mathbb G_n$ on $S(V)$ is compatible    with the
algebra structure $S_{\mathbf{-1}}(V)$ on $S(V)$ $($which is $S_{\mathbf q}(V)$ with $q_{ij}=-1$ for all $i\ne j)$;

\item[$(d)$] $\rho_+(\mathbb C \mathbb G_n)= \rho_-(\mathbb C \mathbb G_n)$, where $\rho_\pm\colon \mathbb C\mathbb G_n \to
\End_\mathbb C S(V)$ are the algebra homomorphisms arising from the actions $\mathop{\triangleright}\nolimits_\pm$.
Moreover, $\rho_-=\rho_+ \circ J$ for some algebra automorphism~$J$ of~$\mathbb C \mathbb G_n$.
\end{enumerate}
\end{Proposition}

\begin{Remark}
\label{rem:injectivity}
By a~powerful result on group actions on integral domains, see Corollary~\ref{cor:premet-dedekind} from the
Appendix taken with $R=\mathbb C$, $A=S(V)$, the action $\mathop{\triangleright}\nolimits_+$ is faithful, and
moreover the corresponding algebra homomorphism $\rho_+\colon \mathbb C \mathbb G_n \to \End_\mathbb C S(V)$ is
injective.
It follows from (d) that $\rho_-\colon \mathbb C \mathbb G_n\to \End_\mathbb C S_{\mathbf{-1}}(V)$ is also injective.
This fact does not readily follow from classical results.
\end{Remark}



At this point, we restrict our attention to f\/inite subgroups~$G$ of $\mathbb G_n$~-- namely, to Shephard--Todd's
imprimitive complex ref\/lection groups $G(m,p,n)$ and the groups $W_{{\mathcal C},{\mathcal C}'}$, introduced
independently in~\cite{BBselecta} and~\cite{KKZ} and def\/ined as follows.
Let $n\ge 1$ and ${\mathcal C}'\subseteq{\mathcal C}$ be two f\/inite subgroups of $\mathbb C^\times$ of orders $\frac
mp$, $m$ respectively.
Then
\begin{gather*}
G(m,p,n)=\{tw\in {\mathcal C}^n\rtimes \mathbb S_n: \det t\in{\mathcal C}'\},
\\
W_{{\mathcal C},{\mathcal C}'} = \{ tw\in {\mathcal C}^n\rtimes \mathbb S_n: \det(tw)\in {\mathcal C}'\}.
\end{gather*}
The similarity of the two def\/initions manifests itself in our f\/irst main result  where a~correspondence, $\mu$,
between the two classes of subgroups of $\mathbb G_n$ is established.

\begin{maintheorem}
\label{thm:main}
Given $n\ge 1$, an even $m\ge 2$ and a~divisor $p$ of $m$, let $G=G(m,p,n)$.
Then there exists a~unique finite subgroup $ \mu(G) \subset \mathbb G_n $ such that the restriction of
$\mathop{\triangleright}\nolimits_-$ onto $\mu(G)$ is mystically equivalent to $\mathop{\triangleright}\nolimits_+$ on~$G$.
In fact,
\begin{gather*}
\mu(G)=W_{\mathcal C,\mathcal C'},
\end{gather*}
where $|{\mathcal C}|=m$, $|{\mathcal C}'|=\frac mp$.
\end{maintheorem}

The def\/inition of the group $\mu(G)$ suggests that the invariants of $\mu(G(m,p,n))=W_{{\mathcal C},{\mathcal C}'}$
should be viewed in the noncommutative algebra $S_{\mathbf{-1}}(V)$, where this group acts via
$\mathop{\triangleright}\nolimits_-$.
To describe these invariants, introduce the following elements in the space $S(V)$:
\begin{gather*}
p_k^{(m)} = \sum\limits_{i=1}^n x_i^{km},
\qquad
k=1,\ldots,n-1,
\qquad
r^{(l)} = x_1^{l}x_2^{l}\cdots x_n^{l}.
\end{gather*}
The classical result of Shephard--Todd and Chevalley asserts that the subalgebra $S(V)^{G(m,p,n)}$ of~$S(V)$ (with
respect to the natural commutative product on~$S(V)$) equals
the polynomial algebra $\mathbb C\big[p_1^{(m)},\ldots,p_{n-1}^{(m)}, r^{(\frac{m}{p})}\big]$.

Our mystic equivalence construction immediately leads to the following result, f\/irst obtained by Kirkman, Kuzmanovich
and Zhang as a~key ingredient in the classif\/ication theorem~\cite[Theorem~1.1]{KKZ}.

\begin{Theorem}
In the notation of Theorem~{\rm \ref{thm:main}}, let $m$ be even.
Then in the algebra $S_{\mathbf{-1}}(V)$,
\begin{enumerate}\itemsep=0pt
\item[$(a)$] the elements $p_1^{(|\mathcal C|)},\ldots,p_{n-1}^{(|\mathcal C|)}, r^{(|\mathcal C'|)}$ are pairwise commuting
invariants of the group $W_{\mathcal C,\mathcal C'}$;

\item[$(b)$] $S_{\mathbf{-1}}(V)^{W_{\mathcal C,\mathcal C'}} = \mathbb C\big[p_1^{(|\mathcal C|)},\ldots,p_{n-1}^{(|\mathcal C|)},r^{(|\mathcal C'|)}\big]$.
\end{enumerate}
\end{Theorem}

\begin{Remark}%\label{rem:EKL}
This result shows that the condition that $m$ is even in Theorem~\ref{thm:main} is important: the symmetric group
$\mathbb S_n=G(1,1,n)$ does not have a~mystical counterpart, and indeed the correspondence $\mu$ cannot be extended to
groups $G(m,p,n)$ where $m$ is odd.
This happens because the space of the invariants of $G(m,p,n)$ in $S(V)$ is not closed under the multiplication in
$S_{\mathbf{-1}}(V)$, and therefore cannot be the space of invariants of a~group acting by automorphisms of
$S_{\mathbf{-1}}(V)$.
\end{Remark}

The groups $G=G(m,p,n)$ and $\mu(G)$ are of the same order $\frac{m^n n!}{p}$, and, informally, they ``look very similar''.
Our next main result makes this informal statement more precise.
We keep the notation from Theorem~\ref{thm:main} and denote by $R$ the ring $\mathbb
Z\big[\frac{1+\mathbf{i}}{2}\big]\subset\mathbb C$.
\begin{maintheorem}
\label{thm:main2}
For all~$G$ as in Theorem~{\rm \ref{thm:main}}, the group rings $RG$ and $R\mu(G)$ are isomorphic.
In particular, $ \mathbb C G \cong \mathbb C \mu(G) $.
\end{maintheorem}
